%%%%%%%%%%%%%%%%%%%%%%%%%%%%%%%%%%%%%%%%%%%%%%%%%%%%%%%%%%%%%%%%%%%%%%%%
%                                                                      %
%     File: Thesis_Contact_Models                                      %
%     Tex Master: Thesis.tex                                           %
%                                                                      %
%     Author: Francisco Castro                                         %
%     Last modified :  2 Jul 2015                                      %
%                                                                      %
%%%%%%%%%%%%%%%%%%%%%%%%%%%%%%%%%%%%%%%%%%%%%%%%%%%%%%%%%%%%%%%%%%%%%%%%
\chapter{Implementation}
\label{chapter:implementation}



%%%%%%%%%%%%%%%%%%%%%%%%%%%%%%%%%%%%%%%%%%%%%%%%%%%%%%%%%%%%%%%%%%%%%%%%
\section{Fast marching}
\label{section:fastMarching}

In order to feed the controller, with a reference path on every problem before its solution is computed, an offline survey of the terrain is conducted in order to build a cost map which takes into account the terrain roughness, the difficulty  mobile robot faces if it attempts at some position and the map inclination. These variables are described below,

Regarding the terrain roughness computation, it is considered that a set of positions around the position of interest play a role on determining the final terrain roughness. If a map $\Map$ is divided into $M$ discretized sets, and if the location where the roughness is computed is said to be $p = (i, j)$, a squared set of size $m$ around $p$ is considered

\begin{equation}
	\bar{\Map}(i, j) = \frac{\sum_{k = i-m, l = j-m}^{i+m, j+m} \Map(k, l)}{m \times m}
\end{equation}



